\documentclass[12pt]{article}

\usepackage[utf8]{inputenc}
\usepackage[T1]{fontenc}
\usepackage{graphicx}
\usepackage{amsmath}
\usepackage{amssymb}
\usepackage{fancyhdr}
\usepackage{geometry}
\usepackage{hyperref}
\usepackage{booktabs}
\usepackage{caption}
\usepackage{titlesec}
\usepackage[expansion=false]{microtype}

\geometry{margin=1in}
\emergencystretch=3em
\sloppy
\raggedbottom
\makeatletter
\setlength{\@fptop}{0pt}
\setlength{\@fpbot}{0pt plus 1fil}
\makeatother

\hypersetup{
    colorlinks=true,
    linkcolor=black,
    citecolor=black,
    urlcolor=black
}

\newcommand{\shorttitle}{Building and Validating the Reddit Kolkata Dataset}

\pagestyle{fancy}
\fancyhf{}
\renewcommand{\headrulewidth}{0pt}
\renewcommand{\footrulewidth}{0.4pt}
\fancyfoot[L]{\shorttitle}
\fancyfoot[R]{Page \thepage}

\titleformat{\section}{\large\bfseries}{\thesection.}{0.6em}{}
\titleformat{\subsection}{\normalsize\bfseries}{\thesubsection.}{0.6em}{}
\titleformat{\subsubsection}{\normalsize\bfseries}{\thesubsubsection.}{0.6em}{}

\newcommand{\sectionbreak}{\clearpage}

\begin{document}

% ---------------------------------------------------------------
% TITLE PAGE
% ---------------------------------------------------------------
\begin{titlepage}
\centering
\vspace*{1cm}
{\large\bfseries MidSem Report Design Project}\\[0.8cm]
{\Large\bfseries Building, Processing, and Validating a Reddit\\
Dataset for r/kolkata and r/KolkataCity: An\\
Incremental Collection and ETL Pipeline}\\[1.5cm]

\includegraphics[width=4cm]{bits_logo.png}\\[1.5cm]

{\bfseries\large Marmik Sapovadia}\\
(2023A7PS0057H)\\[1.5cm]

{\bfseries Under Guidance of}\\
{\bfseries\large Prof. Dipanjan Chakraborty}\\
(BITS Pilani, Hyderabad Campus)\\[1.5cm]

{Course Number: CS F377}
\vfill
\end{titlepage}

% ---------------------------------------------------------------
% CONTENTS
% ---------------------------------------------------------------
{\small\tableofcontents}
\clearpage

% ---------------------------------------------------------------
% ACKNOWLEDGMENTS
% ---------------------------------------------------------------
\section*{ACKNOWLEDGMENTS}
\addcontentsline{toc}{section}{ACKNOWLEDGMENTS}

I would like to thank my project guide, Prof. Dipanjan Chakraborty, for his continual feedback on
the data-processing and validation design decisions described in this report, and for pushing the
project toward a dataset that is verifiably correct rather than merely collected.

I am thankful to the Birla Institute of Technology and Science (BITS), Pilani, Hyderabad Campus,
for the academic environment and computational resources that made this work possible, and to my
peers in the CS F377 course for discussions on deduplication strategy and schema design that
shaped the processing pipeline documented here.

I also acknowledge the maintainers of Arctic Shift, the free archival service that supplied the
raw data this pipeline processes, and the developers of the \texttt{pandas} and \texttt{NumPy}
libraries, on which every processing and validation script in this project depends. Finally, I
acknowledge Baumgartner et al. (2020) and the Pushshift project, whose archival approach Arctic
Shift continues and which this project's own validation methodology draws on.

% ---------------------------------------------------------------
% ABSTRACT
% ---------------------------------------------------------------
\section*{ABSTRACT}
\addcontentsline{toc}{section}{ABSTRACT}

Collecting raw posts and comments from a subreddit is only half of building a usable dataset; the
other half is turning that raw, API-shaped data into clean, de-duplicated, analysis-ready tables
whose correctness can actually be verified. This project is concerned with that second half: the
processing and validation layer of a pipeline that builds a dataset from r/kolkata and
r/KolkataCity, covering 1 April 2025 to the present, on top of raw data fetched primarily from
Arctic Shift, a free, authentication-free, continuously updated Reddit archive, after an initial
batch was collected via BrightData's commercial Datasets API before that source's account quota
was exhausted.

Raw JSONL responses, cached per subreddit and month, are converted into five processed tables —
posts, comments, cross-posts, per-user participation, and per-month activity — through a set of
scripts that re-derive every table from scratch on each run, so that collection gaps can be closed
or new subreddits added without any manual table surgery. A validation step cross-checks these
tables against the collection-state ledger to report duplicate IDs, missing collection periods,
and collection failures, rather than assuming correctness from a successful run. As of this
report, the processed dataset comprises 104,827 posts and 743,126 comments across both
subreddits, 98,855 unique participating users, and zero duplicate post or comment IDs detected
across repeated, independently resumed collection runs.

This report documents the task specification this pipeline was built against, the schema and
deduplication design of the processing layer, the resulting dataset's validation statistics, and
the practical challenges of processing several-hundred-megabyte CSV outputs and reconciling user
activity across two subreddits with very different volumes. This report documents the task
specification the pipeline was built against, the processing and validation methodology, the
resulting dataset statistics, and the engineering challenges specific to correctness-checking a
dataset at this scale.

The remainder of the report is organized as follows. Section~1 introduces why a processing and
validation layer matters as much as collection itself. Section~2 states the aims and scope of this
mid-term phase. Section~3 reviews the literature that informed the schema and validation design.
Section~4 details the processing pipeline, deduplication strategy, and validation methodology.
Section~5 presents the resulting dataset's statistics. Section~6 documents the practical
challenges encountered. Sections~7 and 8 close with future work and concluding remarks, followed
by a supporting appendix and the reference list.

\clearpage

% ---------------------------------------------------------------
% 1. INTRODUCTION
% ---------------------------------------------------------------
\section{Introduction to Dataset Processing and Validation}

\subsection{Raw API Data Is Not a Dataset}
A raw JSONL dump of API responses is not, by itself, something a researcher can run an analysis
against. Reddit's API-shaped JSON carries inconsistent fields across posts, deleted and removed
content masquerading as ordinary rows, and no guarantee that a given post or comment was not
fetched twice across two overlapping collection windows. Turning this raw material into a dataset
that can be trusted requires a processing layer that imposes a fixed schema, removes duplicates
by a stable identifier, and a validation layer that actively checks the result rather than simply
trusting that the collection scripts ran correctly.

\subsection{Why Validation Is a First-Class Concern}
For a multi-month, resumable collection project, correctness bugs are easy to introduce
silently: an interrupted run, a retried request, or a later run with an extended date range can
each introduce duplicate records or leave an undetected gap in coverage if the processing and
validation layer does not actively check for them. Treating validation as an afterthought risks
discovering such problems only when a downstream analysis produces an inexplicable result, long
after the data was collected.

\subsection{This Project}
This project is concerned with the processing and validation layer built on top of the raw data
collected from r/kolkata and r/KolkataCity. It converts cached raw API responses into five
de-duplicated, schema-consistent CSV tables, and produces a validation report that checks those
tables against the project's own collection-state ledger for duplicate IDs, missing periods, and
collection failures, so that the dataset's completeness and correctness claims are backed by an
explicit, re-runnable check rather than an assumption.

% ---------------------------------------------------------------
% 2. AIMS AND SCOPE
% ---------------------------------------------------------------
\section{Aims and Scope}

This mid-term phase is scoped around the processing, derived-table, and validation requirements
of the project's task specification, built on top of raw post and comment data for r/kolkata and
r/KolkataCity collected from 1 April 2025 to the present. The objectives of this mid-term phase
are:

\begin{enumerate}
    \item \textbf{Define a Stable Processed Schema:} Convert raw, API-shaped JSONL into fixed-schema
    posts, comments, and cross-post tables, retaining every metadata field specified in the task
    brief (deletion/removal status, stickied status, score, flair, and so on) without silently
    dropping fields the raw API happens to omit for some records.

    \item \textbf{Deduplicate by Reddit ID:} Ensure that no post or comment appears twice in the
    processed tables even when its collection window was fetched more than once across separate,
    resumed runs, using Reddit's own \texttt{post\_id}/\texttt{comment\_id} as the deduplication
    key.

    \item \textbf{Build Derived Participation Tables:} From the processed posts and comments alone,
    with no additional API calls, construct a per-user participation table (posts/comments per
    subreddit, first-seen/last-seen timestamps) and a per-month activity table, and report user
    overlap between the two subreddits.

    \item \textbf{Validate the Result:} Cross-check the processed tables against the
    collection-state ledger to report, per subreddit, total posts/comments, unique authors,
    earliest/latest post, cross-posts, deleted/removed counts, duplicate IDs, missing collection
    periods, and collection failures.

    \item \textbf{Keep Processing Idempotent:} Ensure every processing and validation script can be
    re-run from scratch at any time, against whatever raw data currently exists on disk, without
    manual intervention or partial-state bugs, so that closing a collection gap or adding a new
    subreddit requires no change to the processing code.
\end{enumerate}

\textbf{Out of scope for this mid-term phase:} any downstream content analysis of the processed
text (sentiment, topic modelling, or similar), reconciling user identities across Reddit-wide
history beyond what this project's own two subreddits contain, and closing the one known
collection gap in the underlying raw data (comments for r/kolkata, June 2025 through December
2025), which is a collection-layer task rather than a processing or validation one and is
discussed further in Section~7.


% ---------------------------------------------------------------
% 3. LITERATURE REVIEW
% ---------------------------------------------------------------
\section{Literature Review and Knowledge Acquisition}

\subsection{Pushshift and the Archive-First Approach to Reddit Data}
The methodological starting point for this project's data source is Baumgartner et al. (2020),
\textit{``The Pushshift Reddit Dataset''}, published at ICWSM, which established the practice of
archiving Reddit content independently of Reddit's own API so that researchers are not subject
to its rate limits or access restrictions. Arctic Shift, the data source this project's raw data
is drawn from, continues this approach after changes to Reddit's API policy in 2023 disrupted
Pushshift's own ingestion.

\textbf{Key Takeaway for This Report:} An archive-first data source shifts the correctness burden
downstream — since the archive itself may contain gaps, inconsistencies, or stale records, the
processing and validation layer built on top of it, which is this report's focus, cannot assume
the raw data is already clean. This motivated treating validation as an explicit, separate step
rather than folding a few sanity checks into the collection scripts themselves.

\subsection{Idempotent, Re-Derivable Processing}
A central design choice in this project's processing layer is that every processed table is
re-derived from the full set of cached raw files on every run, rather than being incrementally
patched. This follows a general data-engineering principle: a pipeline stage that can always be
re-run from its inputs to produce the same output is easier to reason about and debug than one
that accumulates state through incremental patches, because a bug in an incremental update can
leave the output permanently inconsistent with its inputs in a way that is hard to detect later.

\subsection{Deduplication by Stable Identifier}
Using Reddit's own \texttt{post\_id}/\texttt{comment\_id} as the deduplication key, rather than a
derived or positional key, follows directly from these identifiers being globally stable and
assigned by Reddit itself at creation time — any two fetches of the same post, however far apart
in time or however overlapping their collection windows, return the same ID. This is the same
identifier-stability property that made Pushshift-style archives usable as a merge target for
data collected across multiple, independently-run sessions in the first place.


% ---------------------------------------------------------------
% 4. METHODOLOGY
% ---------------------------------------------------------------
\section{Processing and Validation Methodology}

\subsection{Raw Data as the Single Source of Truth}
Every processed table in this project is derived exclusively from the raw JSONL files cached
under \texttt{raw/<subreddit>/<month>/} during collection. Collection itself used two sources
over the project's lifetime: BrightData's Datasets API, which supplied the project's earliest
raw data until the account's request quota was exhausted, after which collection moved to Arctic
Shift, a free, authentication-free archive with no such quota, which supplied the remainder and
is the source of essentially all data in the processed tables (the BrightData path is retained
in the collection code but unused in practice). No processing or validation script issues a
network request; all of them operate purely on what is already on disk, which keeps a full
re-processing run fast and makes the raw cache the single, authoritative source of truth for
every downstream table.

\subsection{Posts and Cross-Posts}
\texttt{process\_posts.py} reads every month's raw posts file for both subreddits and emits
\texttt{posts.csv} with a fixed column schema — subreddit, post ID, timestamp, author, title, post
text, flair, post type, score, upvote ratio, comment count, permalink, external URL/domain,
cross-post status, stickied status, and deletion/removal status — filling in \texttt{null} rather
than omitting a column when a given raw record lacks a field. Every post flagged as a cross-post
additionally emits a row in \texttt{crossposts.csv}, carrying the original post's ID, subreddit,
author, URL, and date where available, without triggering any additional collection of the
originating subreddit.

\subsection{Comments}
\texttt{process\_comments.py} reads the two-level comment tree already reconstructed during
collection (top-level comments plus their direct replies; deeper reply chains are not collected,
as documented in the project README) and flattens it into \texttt{comments.csv}, one row per
comment, with \texttt{depth} values of 0 or 1 and \texttt{parent\_id} preserved so that a
depth-1 reply tree can be reconstructed directly from the processed table without re-parsing raw
data.

\subsection{Deduplication by Reddit ID}
Both \texttt{process\_posts.py} and \texttt{process\_comments.py} deduplicate by Reddit's own
\texttt{post\_id}/\texttt{comment\_id} while appending rows, so that a post or comment collected
twice (because two collection runs' date windows overlapped, or because a crash caused a partial
re-fetch) is written to the processed table exactly once, regardless of how many times it appears
across the raw JSONL cache.

\subsection{Derived Tables: Users and Monthly Activity}
\texttt{build\_users.py} scans the processed posts and comments tables once to build
\texttt{users.csv}: one row per username, with per-subreddit post and comment counts and
first-seen/last-seen timestamps, plus a console summary of users active in both subreddits,
only r/kolkata, and only r/KolkataCity. \texttt{build\_monthly\_activity.py} performs the
analogous aggregation by calendar month rather than by user, producing
\texttt{monthly\_activity.csv} with posts, comments, unique active users, and newly-first-seen
users per subreddit per month. Both scripts operate purely on the processed tables, not the raw
cache, and are re-run in full whenever the underlying posts or comments table changes.

\subsection{Validation Against the Collection-State Ledger}
\texttt{validate.py} is the correctness check for the whole pipeline: it re-scans
\texttt{collection\_state.csv} (the per-\texttt{(subreddit, data\_type, month)} ledger maintained
during collection) alongside the processed tables, and reports, per subreddit, the headline
statistics in Section~5, together with any duplicate post/comment IDs found, any month flagged as
other than \texttt{complete} in the state ledger (rendered as a human-readable missing-period
string), and every row in \texttt{collection\_state.csv} with status \texttt{failed} or
\texttt{partial}, recorded verbatim as a collection failure. This step is what turns "the
collection scripts ran without crashing" into an explicit, checkable claim about what was and was
not actually collected.


% ---------------------------------------------------------------
% 5. RESULTS
% ---------------------------------------------------------------
\section{Dataset and Validation Results}

\subsection{Processed Table Sizes and Headline Statistics}
Table~\ref{tab:validation} reports the validation statistics produced by \texttt{validate.py} for
both subreddits (generated 1 October 2026). The processed comment table is approximately 226~MB
and the post table approximately 74~MB; across both subreddits, the pipeline has produced 104,827
processed post rows and 743,126 processed comment rows, with zero duplicate IDs detected in
either table.

\begin{table}[h]
\centering
\small
\caption{Validation statistics for both subreddits (generated 2026-10-01).}
\label{tab:validation}
\begin{tabular}{lrr}
\toprule
\textbf{Metric} & \textbf{r/kolkata} & \textbf{r/KolkataCity} \\
\midrule
Total posts & 91{,}328 & 13{,}499 \\
Total comments & 566{,}456 & 176{,}670 \\
Unique authors & 71{,}883 & 38{,}808 \\
Earliest post & 2025-04-01 & 2025-04-14 \\
Latest post & 2026-09-26 & 2026-09-25 \\
Cross-posts & 0 & 93 \\
Posts with external URLs & 24 & 0 \\
Deleted/removed posts & 12{,}577 & 1{,}073 \\
Deleted/removed comments & 38{,}945 & 19{,}820 \\
Duplicate post/comment IDs & 0 & 0 \\
\bottomrule
\end{tabular}
\end{table}
\vspace{-0.5cm}

\subsection{User Participation Across Subreddits}
\texttt{build\_users.py} produced 98,855 unique usernames across both subreddits. Of these,
11,836 users post or comment in both r/kolkata and r/KolkataCity, 60,047 are active only in
r/kolkata, and 26,972 are active only in r/KolkataCity. Relative to each subreddit's own author
count from Table~\ref{tab:validation} (71,883 for r/kolkata, 38,808 for r/KolkataCity), the
cross-subreddit overlap represents roughly 16\% of r/kolkata's author base and roughly 30\% of
r/KolkataCity's, consistent with r/KolkataCity functioning as a smaller, more tightly overlapping
community relative to the larger, more general-purpose r/kolkata.

\subsection{Discussion}
The zero-duplicate-ID result in Table~\ref{tab:validation}, reproduced consistently by
\texttt{validate.py} across repeated runs on top of a raw cache that was itself built from many
separate, independently resumed collection sessions, is the main correctness evidence this
processing layer provides: it confirms that the deduplication-by-Reddit-ID strategy
(Section~4.4) holds in practice, not merely in design, under exactly the kind of overlapping,
interrupted collection history this dataset actually has.


% ---------------------------------------------------------------
% 6. CHALLENGES
% ---------------------------------------------------------------
\section{Challenges and Strategic Optimizations}

\subsection{Reconciling Inconsistent Raw Schemas}
Raw records collected via Arctic Shift do not all carry the same set of fields — a field such as
\texttt{upvote\_ratio} or \texttt{stickied} may be present on a post fetched recently but absent
on one fetched earlier in the project, before that field was wired into the collection client, or
on the small number of posts originally sourced via the now-unused BrightData path, before that
source's account quota was reached and collection moved to Arctic Shift. Rather than
dropping a column whenever any record lacked it, \texttt{process\_posts.py} fills missing fields
with \texttt{null} on a per-record basis, so the processed schema stays fixed across the whole
table and a consumer of \texttt{posts.csv} can distinguish "this field is genuinely absent for
this record" from "this column does not exist."

\subsection{Deduplicating Across Independently Resumed Runs}
Because collection for this project happened across many separate sessions, often resumed after
an interruption or extended with a later end date, the raw JSONL cache contains a meaningful
number of posts and comments fetched more than once. Verifying that deduplication actually
collapses these correctly — rather than merely trusting the dedup logic — required deliberately
cross-checking \texttt{validate.py}'s duplicate-ID count against spot checks of the raw cache for
known-overlapping collection windows, since a silent dedup failure would not have caused any
visible error, only an inflated and incorrect processed table.

\subsection{Processing Very Large Tables Without Loading Them Fully Into Memory}
With the comment table alone around 226~MB, early versions of the processing scripts that
materialized the full table in memory before writing it out were noticeably slower and more
memory-hungry than necessary. Restructuring both \texttt{process\_posts.py} and
\texttt{process\_comments.py} to append rows to disk incrementally as each raw file is read,
rather than building one large in-memory structure first, kept processing within the memory
budget of the development machine and made re-processing after a collection gap is closed
considerably faster.

\subsection{Interpreting a Validation Report That Flags Expected Gaps}
\texttt{validate.py} reports every non-\texttt{complete} row in the collection-state ledger as a
missing period or collection failure, including the nine months of r/kolkata comments that were
deliberately never attempted as a scope decision (Section~6.3 of the companion collection-focused
report) rather than collected and lost. Distinguishing a genuinely unexpected gap from an
intentionally deferred one required cross-referencing the validation report against the project's
own scope decisions each time it was regenerated, rather than treating every reported gap as a
bug to fix immediately.


% ---------------------------------------------------------------
% 7. FUTURE WORK
% ---------------------------------------------------------------
\section{Future Work}

As the project advances beyond the mid-term review, the following directions are planned for the
processing and validation layer specifically:

\begin{itemize}
    \item \textbf{Automated Gap Classification:} Extend \texttt{validate.py} to distinguish
    deliberately deferred collection gaps from genuinely unexpected ones automatically, rather
    than requiring manual cross-referencing against scope decisions each time the report is
    regenerated.

    \item \textbf{Schema Versioning:} Introduce an explicit schema version field in the processed
    tables, so that a future change to the set of retained columns does not silently change the
    meaning of an already-published \texttt{posts.csv} or \texttt{comments.csv}.

    \item \textbf{Cross-Validation Against a Second Deduplication Pass:} Implement an independent,
    brute-force deduplication check over the full raw cache as a periodic audit, to catch any
    edge case the streaming deduplication logic in \texttt{process\_posts.py}/
    \texttt{process\_comments.py} might miss.

    \item \textbf{Processing-Time Benchmarking:} Measure and report the wall-clock cost of a full
    re-processing run as the raw cache grows, to anticipate when the current single-pass,
    streaming-write design will need to be replaced by a more incremental approach.

    \item \textbf{Extending Derived Tables to New Subreddits:} Once the pipeline is extended to
    additional subreddits (Section~7 of the companion collection-focused report), verify that
    \texttt{build\_users.py} and \texttt{build\_monthly\_activity.py} correctly extend their
    output schema (new \texttt{posts\_<subreddit>}/\texttt{comments\_<subreddit>} columns) without
    requiring any code change, as designed.
\end{itemize}


% ---------------------------------------------------------------
% 8. CONCLUSION
% ---------------------------------------------------------------
\section{Conclusion}

This mid-term phase built and validated the processing layer that turns raw, API-shaped Reddit
data for r/kolkata and r/KolkataCity into a clean, de-duplicated dataset: five processed tables
derived entirely from cached raw data with no re-querying of any API, and a validation step that
checks the result against the project's own collection-state ledger rather than assuming
correctness. As of this report, the processed dataset comprises 104,827 posts and 743,126
comments across both subreddits, 98,855 unique participating users with a documented overlap of
11,836 users active in both communities, and zero duplicate post or comment IDs detected despite
a long, repeatedly interrupted and resumed collection history.

The central engineering result of this phase is that the deduplication-by-Reddit-ID strategy and
the re-derive-from-scratch processing design (Section~4) hold up under exactly the conditions
they were designed for — overlapping, independently resumed collection windows — as confirmed
by \texttt{validate.py}'s zero-duplicate finding rather than assumed from the pipeline's design
alone. The immediate next steps are to make the validation report distinguish deliberate
collection-scope decisions from genuine gaps automatically, and to extend the same processing and
validation design to additional subreddits as the collection layer grows, as outlined in
Section~7.

Beyond the specific numbers, this phase demonstrates that treating data validation as a first-class,
independently re-runnable pipeline stage, rather than an assumption built into the collection
scripts, is what makes it possible to state this dataset's completeness and correctness as a
checked fact rather than a hope — a distinction that matters considerably more once the dataset
is handed off to downstream analysis that will not re-verify these properties itself.


% ---------------------------------------------------------------
% 9. APPENDIX
% ---------------------------------------------------------------
\section{Appendix}

\subsection*{A.1 Processed Table Schema Summary}
Table~\ref{tab:schema} summarizes the five processed tables produced by this pipeline and the
script responsible for each.

\begin{table}[h]
\centering
\small
\caption{Processed tables and their generating scripts.}
\label{tab:schema}
\begin{tabular}{@{}llp{6.3cm}@{}}
\toprule
\textbf{Table} & \textbf{Script} & \textbf{Key fields} \\
\midrule
\texttt{posts.csv} & \texttt{process\_posts.py} & subreddit, post\_id, timestamp, author, title,
score, flair, permalink, deletion/removal status \\
\texttt{crossposts.csv} & \texttt{process\_posts.py} & post\_id, original\_post\_id,
original\_subreddit, original\_author, original\_url \\
\texttt{comments.csv} & \texttt{process\_comments.py} & subreddit, post\_id, comment\_id,
parent\_id, depth (0/1), author, score \\
\texttt{users.csv} & \texttt{build\_users.py} & username, per-subreddit post/comment counts,
first\_seen, last\_seen \\
\texttt{monthly\_activity.csv} & \texttt{build\_monthly\_activity.py} & month, subreddit, posts,
comments, unique\_active\_users, new\_users \\
\bottomrule
\end{tabular}
\end{table}

\subsection*{A.2 User Overlap Between Subreddits}
Table~\ref{tab:overlap} reports the user-overlap breakdown computed from \texttt{users.csv}.

\begin{table}[h]
\centering
\small
\caption{User overlap between r/kolkata and r/KolkataCity.}
\label{tab:overlap}
\begin{tabular}{lr}
\toprule
\textbf{Category} & \textbf{Users} \\
\midrule
Total unique users (either subreddit) & 98{,}855 \\
Users active in both subreddits & 11{,}836 \\
Users active only in r/kolkata & 60{,}047 \\
Users active only in r/KolkataCity & 26{,}972 \\
\bottomrule
\end{tabular}
\end{table}

\subsection*{A.3 Processing Environment}
Table~\ref{tab:env} summarizes the software environment used to produce and validate the
processed dataset reported above.

\begin{table}[h]
\centering
\caption{Software environment for the processing and validation layer.}
\label{tab:env}
\begin{tabular}{@{}lp{8.5cm}@{}}
\toprule
\textbf{Setting} & \textbf{Value} \\
\midrule
Raw data source consumed & Arctic Shift JSONL cache (\texttt{raw/<subreddit>/<month>/}) \\
Processed comment table size & $\approx$226~MB \\
Processed post table size & $\approx$74~MB \\
Deduplication key & Reddit \texttt{post\_id} / \texttt{comment\_id} \\
Comment tree depth retained & 2 levels (top-level + direct reply) \\
Validation cross-check source & \texttt{collection\_state.csv} \\
Core libraries & Python 3, \texttt{pandas}, \texttt{NumPy} \\
\bottomrule
\end{tabular}
\end{table}

% ---------------------------------------------------------------
% 10. REFERENCES
% ---------------------------------------------------------------
\section{References}

\begin{enumerate}
    \item Baumgartner, J., Zannettou, S., Keegan, B., Squire, M., \& Blackburn, J. (2020).
    \textit{The Pushshift Reddit Dataset}. Proceedings of the International AAAI Conference on
    Web and Social Media (ICWSM), 14, 830--839.

    \item Heitmann, A. \textit{Arctic Shift: A Continuously Updated Archive of Reddit Content}.
    GitHub repository. \url{https://github.com/ArthurHeitmann/arctic_shift}

    \item Reddit, Inc. \textit{Reddit API Documentation}. \url{https://www.reddit.com/dev/api/}

    \item McKinney, W. (2010). \textit{Data Structures for Statistical Computing in Python}.
    Proceedings of the 9th Python in Science Conference (SciPy), 56--61.

    \item Harris, C. R., Millman, K. J., van der Walt, S. J., Gommers, R., Virtanen, P.,
    Cournapeau, D., et al. (2020). \textit{Array Programming with NumPy}. Nature, 585,
    357--362.

    \item Kleppmann, M. (2017). \textit{Designing Data-Intensive Applications}. O'Reilly Media.
    (Idempotent, re-derivable pipeline design reference for Section~3.2 and Section~4.1.)
\end{enumerate}

\end{document}
